\section*{Capítulo \ref{ch:b1:integration} --- Integración en un segmento}

\begin{solution}{exo:b1:integration:1}
$\dfrac{1}{x^2-4} = \dfrac{1/4}{x - 2} - \dfrac{1/4}{x+2}$
(por tapado), luego
\[
\int_0^1 \frac{\dd x}{x^2 - 4}
= \frac14\Bigl[\ln\abs{x-2} - \ln\abs{x+2}\Bigr]_0^1
= \frac14\Bigl(\ln\frac{1}{3} - \ln 1\Bigr) = -\frac{\ln 3}{4}.
\]

Por partes ($u' = 1$, $v = \ln t$): $\int_1^{\eu} \ln t\,\dd t =
\bigl[t\ln t\bigr]_1^{\eu} - \int_1^{\eu} \dd t = \eu - (\eu - 1) =
1$.

Por partes ($u' = \sin t$, $v = t$): $\int_0^\pi t\sin t\,\dd t =
\bigl[-t\cos t\bigr]_0^\pi + \int_0^\pi \cos t\,\dd t = \pi + 0 =
\pi$.
\end{solution}

\begin{solution}{exo:b1:integration:2}
\omterm{thm:b1:integration:parts}{Cambio de variable} $u = t^2 + 1$, $\dd u = 2t\,\dd t$:
\[
\int_0^1 \frac{t\,\dd t}{(t^2+1)^2}
= \frac12 \int_1^2 \frac{\dd u}{u^2}
= \frac12\Bigl[-\frac1u\Bigr]_1^2 = \frac14 .
\]
Con $u = \sin t$: $\int_0^{\pi/2} \cos^3 t\,\dd t = \int_0^{\pi/2}
(1 - \sin^2 t)\cos t\,\dd t = \bigl[\sin t - \frac{\sin^3
t}{3}\bigr]_0^{\pi/2} = 1 - \frac13 = \frac23$.
\end{solution}

\begin{solution}{exo:b1:integration:3}
$u_n = \dfrac1n \sum_{k=1}^{n} \dfrac{1}{1 + (k/n)^2}$: una suma de
Riemann de $x \mapsto \frac{1}{1+x^2}$ en $\intcc{0}{1}$, luego
$u_n \to \int_0^1 \frac{\dd x}{1+x^2} = \arctan 1 = \dfrac\pi4$.

$\ln v_n = \dfrac1n \sum_{k=1}^{n} \ln\dfrac{n+k}{n} = \dfrac 1n
\sum_{k=1}^n \ln\Bigl(1 + \dfrac kn\Bigr) \to \int_0^1 \ln(1+x)\,\dd
x = \bigl[(1+x)\ln(1+x) - x\bigr]_0^1 = 2\ln 2 - 1$. Por tanto,
$v_n \to \eu^{2\ln 2 - 1} = \dfrac 4\eu$. (Comprobación de la
identificación: $\frac{(2n)!}{n!\,n^n} = \prod_{k=1}^{n} \frac{n+k}{n}$.)
\end{solution}

\begin{solution}{exo:b1:integration:4}
El límite es $0$. Sean $M = \sup \abs f$ y
$\varepsilon \in \intoo{0}{1}$. Córtese en $1 - \varepsilon$:
\[
\Bigl| \int_0^1 x^n f \Bigr|
\leq \int_0^{1 - \varepsilon} x^n \abs f + \int_{1-\varepsilon}^1
x^n \abs f
\leq M\,(1-\varepsilon)^n + M\varepsilon .
\]
Como $(1 - \varepsilon)^n \to 0$ (\cref{exo:b1:seq:3}), el límite
superior del miembro izquierdo es $\leq M\varepsilon$ para todo
$\varepsilon$: la \omterm{thm:b1:integration:def}{integral} tiende a $0$.
\end{solution}

\begin{solution}{exo:b1:integration:5}
$Q(\lambda) = \int_a^b (f + \lambda g)^2 = \int f^2 + 2\lambda \int
fg + \lambda^2 \int g^2 \geq 0$ para todo $\lambda$. Si
$\int g^2 = 0$, entonces $g = 0$ (positividad estricta,
\cref{thm:b1:integration:props}~(4)) y la desigualdad es
$0 \leq 0$. En caso contrario, $Q$ es un verdadero trinomio de
segundo grado, siempre $\geq 0$: su discriminante es $\leq 0$, es
decir, $\bigl(\int fg\bigr)^2 \leq \int f^2 \int g^2$.

Hay igualdad si y solo si se anula el discriminante, si y solo si
$Q(\lambda_0) = 0$ para algún $\lambda_0$, es decir,
$\int (f + \lambda_0 g)^2 = 0$, es decir (otra vez la positividad
estricta), $f = -\lambda_0 g$: hay igualdad exactamente cuando $f$ y
$g$ son proporcionales.
\end{solution}

\begin{solution}{exo:b1:integration:6}
Sea $\Phi(a) = \int_a^{a+T} f$. Por el teorema fundamental
(\cref{thm:b1:integration:ftc}), $\Phi$ es \omterm{def:b1:derivative:def}{derivable} con
$\Phi'(a) = f(a + T) - f(a) = 0$: constante.

Para $x > 0$, escríbase $x = nT + r$, $0 \leq r < T$
($n = \lfloor x/T \rfloor$). Chasles:
\[
\int_0^x f = n \int_0^T f + \int_{nT}^{nT + r} f,
\qquad
\Bigl| \int_{nT}^{nT+r} f \Bigr| \leq T \sup_{\intcc{0}{T}} \abs f
= C .
\]
Entonces $\frac 1x \int_0^x f = \frac{nT}{x}\cdot\frac 1T \int_0^T f +
O\bigl(\frac 1x\bigr)$, y $\frac{nT}{x} \to 1$: el límite es
$\frac1T \int_0^T f$.
\end{solution}

\begin{solution}{exo:b1:integration:7}
Sea $g(x) = f(x) - x$: \omterm{def:b1:continuity:continuous}{continua}, con
\[
\int_0^1 g = \int_0^1 f - \frac12 = 0 .
\]
Si $g$ no se anulara nunca, el teorema del valor intermedio forzaría
un signo constante (una función \omterm{def:b1:continuity:continuous}{continua} en un \omterm{prop:b1:reals:intervals}{intervalo} que toma
los dos signos se anula); digamos $g > 0$. Entonces, por la
positividad estricta (\cref{thm:b1:integration:props}~(4) aplicada a
$g > 0$, que da $\int g > 0$): contradicción con $\int g = 0$. Luego
$g(c) = 0$ para algún $c$: $f(c) = c$.
\end{solution}

\begin{solution}{exo:b1:integration:8}
\begin{enumerate}
  \item Por partes con $u' = \sin t$, $v = \sin^{n-1} t$:
        \[
        W_n = \bigl[-\cos t\sin^{n-1}t\bigr]_0^{\pi/2}
        + (n-1)\int_0^{\pi/2} \cos^2 t\,\sin^{n-2} t\,\dd t
        = (n-1)(W_{n-2} - W_n),
        \]
        luego $nW_n = (n-1)W_{n-2}$. De $W_0 = \frac\pi2$, $W_1 = 1$:
        \[
        W_{2p} = \frac{(2p-1)(2p-3)\cdots 1}{(2p)(2p-2)\cdots 2}\,
        \frac{\pi}{2} = \frac{(2p)!}{4^p (p!)^2}\,\frac\pi2,
        \qquad
        W_{2p+1} = \frac{(2p)(2p-2)\cdots 2}{(2p+1)(2p-1)\cdots 3}
        = \frac{4^p (p!)^2}{(2p+1)!} .
        \]
  \item En $\intoo{0}{\frac\pi2}$, $0 < \sin t < 1$, luego
        $\sin^{n+1} < \sin^n$ y $(W_n)$ es (estrictamente)
        decreciente y positiva. Encuadrando con la recurrencia:
        \[
        \frac{n}{n+1} = \frac{W_{n+1}}{W_{n-1}}
        \leq \frac{W_{n+1}}{W_n} \leq 1
        \quad\implies\quad \frac{W_{n+1}}{W_n} \to 1 .
        \]
        Invariante: $a_n = (n+1)W_{n+1}W_n$ cumple $a_n = a_{n-1}$
        por la recurrencia $(n+1)W_{n+1} = nW_{n-1}$, luego
        $a_n = a_0 = 1 \cdot W_1 W_0 = \frac\pi2$. Entonces
        \[
        n W_n^2 \sim (n+1) W_{n+1} W_n = \frac\pi2
        \quad\implies\quad
        W_n \sim \sqrt{\frac{\pi}{2n}} .
        \]
\end{enumerate}
\end{solution}

\begin{solution}{exo:b1:integration:9}
\begin{enumerate}
  \item En $\intoo{0}{\pi}$: $x > 0$, $a - bx = b(\frac ab - x) =
        b(\pi - x) > 0$ y $\sin x > 0$, luego el integrando es $> 0$
        y $I_n > 0$ (positividad estricta). Cota: en
        $\intcc{0}{\pi}$, $x \leq \pi$ y $a - bx \leq a$, luego
        $P \leq \frac{\pi^n a^n}{n!}$ y
        $I_n \leq \pi\,\frac{(\pi a)^n}{n!}$, que tiende a $0$ (el
        factorial gana al término geométrico: es el término general
        de la serie exponencial convergente, cf.\
        \cref{ex:b1:seq:e}); en particular, $I_n < 1$ para $n$ grande.
  \item Desarróllese $x^n(a - bx)^n = \sum_{j=0}^{n} \binom nj
        a^{n-j} (-b)^j x^{n+j}$: así, $P = \frac{1}{n!}\sum_j c_j
        x^{n+j}$ con $c_j$ enteros. Entonces $P^{(k)}(0) = 0$ para
        $k < n$ (valuación) y, para $n \leq k \leq 2n$,
        $P^{(k)}(0) = \frac{k!}{n!} c_{k-n}$, un entero porque
        $n! \mid k!$. Además, $P(\pi - x) = P(x)$ (sustitúyase:
        $\pi - x$ intercambia los factores, usando
        $a - b(\pi - x) = bx$), luego $P^{(k)}(\pi) = \pm
        P^{(k)}(0)$: también enteros.
  \item Con $Q = P - P'' + P^{(4)} - \dots$ (finita: $P$ tiene
        grado $2n$): $Q + Q'' = P$, y
        \[
        \bigl(Q'\sin x - Q\cos x\bigr)' = (Q + Q'')\sin x = P\sin x .
        \]
        Por tanto, $I_n = \bigl[Q'\sin x - Q\cos x\bigr]_0^{\pi} =
        Q(\pi) + Q(0)$, una suma de valores $P^{(2k)}$ en $0$ y en
        $\pi$: un entero por (2).
  \item Para $n$ grande, $I_n$ es un entero con $0 < I_n < 1$:
        imposible. La hipótesis $\pi = \frac ab$ falla: $\pi$ es
        irracional.
\end{enumerate}
\end{solution}

\begin{solution}{exo:b1:integration:10}
Caso $C^1$: por partes,
\[
\int_a^b f(t)\sin\lambda t\,\dd t
= \Bigl[-f(t)\frac{\cos\lambda t}{\lambda}\Bigr]_a^b
+ \frac{1}{\lambda}\int_a^b f'(t)\cos\lambda t\,\dd t,
\]
acotado en valor absoluto por
$\frac{2\sup\abs f + (b - a)\sup\abs{f'}} {\lambda} \to 0$.

Caso \omterm{def:b1:continuity:continuous}{continuo}: sea $\varepsilon > 0$ y tómese una \omterm{def:b1:integration:step}{función
escalonada} $\varphi$ con $\abs{f - \varphi} \leq \varepsilon$
(el~\cref{thm:b1:integration:approx} proporciona
$\varphi \leq f \leq \psi$ con hueco $\leq\varepsilon$; tómese
$\varphi$). Entonces
\[
\Bigl|\int_a^b f\sin\lambda t\Bigr|
\leq \int_a^b \abs{f - \varphi}
+ \Bigl|\int_a^b \varphi \sin\lambda t\Bigr|
\leq (b-a)\varepsilon
+ \sum_i \abs{c_i}\,\Bigl|\int_{x_{i-1}}^{x_i} \sin\lambda t\,\dd
t\Bigr| ,
\]
y cada $\bigl|\int \sin \lambda t\bigr| = \bigl|\frac{\cos\lambda
x_{i-1} - \cos\lambda x_i}{\lambda}\bigr| \leq \frac{2}{\lambda}$:
el segundo término tiende a $0$. Luego el límite superior es
$\leq (b-a)\varepsilon$ para todo $\varepsilon$: el límite es $0$.
\end{solution}

\begin{solution}{exo:b1:integration:11}
Sean $m = \min f$ y $M = \max f$, alcanzados por el teorema de los
valores extremos. Como $g \geq 0$: $m\,g \leq fg \leq M\,g$, luego,
por monotonía,
\[
m \int_a^b g \;\leq\; \int_a^b fg \;\leq\; M \int_a^b g .
\]
Si $\int_a^b g = 0$: la positividad estricta
(\cref{thm:b1:integration:props}~(4)) obliga a $g \equiv 0$, los dos
miembros se anulan y sirve cualquier $c$. En caso contrario,
$t = \frac{\int fg}{\int g}$ está en
$\intcc{m}{M} = f(\intcc{a}{b})$
(\cref{thm:b1:continuity:evt,thm:b1:continuity:ivt}), luego
$t = f(c)$ para algún $c$.

El signo importa: en $\intcc{-1}{1}$ con $f(t) = g(t) = t$:
$\int fg = \int_{-1}^1 t^2 = \frac23$, mientras que
$f(c)\int_{-1}^1 t\,\dd t = 0$ para todo $c$.
\end{solution}

\begin{solution}{exo:b1:integration:12}
Si $f \equiv 0$, la afirmación es vacía (todo punto es un cero). Así
que supóngase $f \not\equiv 0$ y que tiene a lo sumo $n$ ceros
distintos en $\intoo{a}{b}$; sean $z_1 < \dots < z_m$ ($m \leq n$)
aquellos ceros en los que $f$ \emph{cambia de signo} (posiblemente
ninguno). Póngase $P(t) = \prod_{i=1}^{m}(t - z_i)$ (producto vacío
$= 1$), de grado $m \leq n$. En cada subintervalo determinado por los
$z_i$, tanto $f$ como $P$ tienen signo constante, y los dos cambian
de signo al cruzar algún $z_i$: el producto $fP$ tiene un signo
constante en todo $\intoo{a}{b}$. Siendo \omterm{def:b1:continuity:continuous}{continuo}, no idénticamente
nulo y de signo constante, cumple $\bigl|\int_a^b fP\bigr| > 0$ (la
positividad estricta aplicada a $\abs{fP}$). Pero $\int f P$ es una
combinación lineal de los momentos $\int f\,t^k$, $k \leq n$, todos
nulos: contradicción. Luego $f$ tiene al menos $n + 1$ ceros
distintos en $\intoo{a}{b}$.
\end{solution}

\begin{solution}{pb:b1:integration:1}
\textbf{1.} Sea $N = \lceil 2c \rceil$. Para $n \geq N$:
$\frac{c^{n+1}/(n+1)!}{c^n/n!} = \frac{c}{n+1} \leq \frac12$, luego
$0 < \frac{c^n}{n!} \leq \frac{c^N}{N!}\,2^{-(n - N)} \to 0$: por
emparedado.

\textbf{2.} Fíjese $k$; inducción sobre $m$. Para $m = 0$:
$\int_0^1 x^k = \frac{1}{k+1} = \frac{k!\,0!}{(k+1)!}$. Paso, por
partes ($u = (1-x)^m$, $v' = x^k$):
\[
\int_0^1 x^k (1-x)^m \dd x
= \Bigl[\frac{x^{k+1}}{k+1}(1-x)^m\Bigr]_0^1
+ \frac{m}{k+1}\int_0^1 x^{k+1}(1-x)^{m-1}\dd x
= \frac{m}{k+1}\cdot\frac{(k+1)!\,(m-1)!}{(k+m+1)!} ,
\]
que es $\frac{k!\,m!}{(k+m+1)!}$.

\textbf{3.} $k = m = n$:
$\int_0^1 (x(1-x))^n = \frac{(n!)^2}{(2n+1)!} =
\frac{1}{(2n+1)\binom{2n}{n}}$. Como $x(1-x) \leq \frac14$ en
$\intcc{0}{1}$, la \omterm{thm:b1:integration:def}{integral} es $\leq 4^{-n}$, de donde
$\binom{2n}{n} \geq \frac{4^n}{2n+1}$ --- acorde con
$\binom{2n}{n}^{1/n} \to 4$ (\cref{pb:b1:seq:1}).

\textbf{4.} El integrando es \omterm{def:b1:continuity:continuous}{continuo}, $\geq 0$ y positivo en
$\intoo{0}{1}$, luego no es idénticamente nulo: su \omterm{thm:b1:integration:def}{integral} es $> 0$
(\cref{thm:b1:integration:props}~(4)). \omterm{def:b1:reals:bounds}{Cota superior}:
$(x(1-x))^n \leq 4^{-n}$ y $g \leq \sup\abs g$, y después la
monotonía.

\textbf{5.} $A_0 = \eu - 1$;
$A_1 = [x\eu^x]_0^1 - \int_0^1 \eu^x = \eu - (\eu - 1) = 1$. Por
partes: $A_n = [x^n \eu^x]_0^1 - n\int_0^1 x^{n-1}\eu^x =
\eu - n A_{n-1}$. Cotas: el integrando es positivo, luego $A_n > 0$;
y $\eu^x \leq \eu$ da $A_n \leq \eu\int_0^1 x^n = \frac{\eu}{n+1}$.

\textbf{6.} $A_0 = -1 + 1\cdot\eu$. Si
$A_{n-1} = \alpha_{n-1} + \beta_{n-1}\eu$ con entradas enteras,
entonces
\[
A_n = \eu - n\alpha_{n-1} - n\beta_{n-1}\eu
= \underbrace{(-n\,\alpha_{n-1})}_{\alpha_n}
+ \underbrace{(1 - n\,\beta_{n-1})}_{\beta_n}\,\eu ,
\]
los dos enteros.

\textbf{7.} Si $\eu = \frac pq$: $q A_n = q\alpha_n + p\beta_n
\in \Z$, y $0 < qA_n \leq \frac{q\eu}{n+1} < 1$ para $n$ grande: un
entero estrictamente entre $0$ y $1$ --- imposible. Luego
$\eu \notin \Q$. En el~\cref{exo:b1:seq:9} el entero atrapado era
$q!\,\frac pq - q!\,a_q$; aquí es $qA_n$: la misma trampa, con
combustible \omterm{thm:b1:integration:def}{integral}.

\textbf{8.} $n = 0$: $R_0 = \int_0^1 \eu^t = \eu - 1$, luego
$\eu = 1 + R_0$. Por partes ($u = \eu^t$, $v = -\frac{(1-t)^{n+1}}{n+1}$):
\[
R_n = \frac{1}{n!}\Bigl(\Bigl[-\frac{(1-t)^{n+1}}{n+1}
\eu^t\Bigr]_0^1 + \frac{1}{n+1}\int_0^1 (1-t)^{n+1}\eu^t\Bigr)
= \frac{1}{(n+1)!} + R_{n+1} ,
\]
de modo que la fórmula se propaga de $n$ a $n + 1$. Cotas:
$1 \leq \eu^t \leq \eu$ en $\intcc{0}{1}$ y
$\int_0^1 (1-t)^n = \frac{1}{n+1}$ dan
$\frac{1}{(n+1)!} \leq R_n \leq \frac{\eu}{(n+1)!}$.

\textbf{9.} $\sum_{k=0}^{9} \frac{1}{k!} = \frac{986410}{362880}
= 2.71828152\dots$, y
\[
\frac{1}{10!} = 2.76\cdot10^{-7} \leq R_9 \leq
\frac{2.75}{10!} = 7.58\cdot10^{-7} ,
\]
luego $2.7182818 \leq \eu \leq 2.7182823$: con demostración,
$\eu = 2.718282$ con seis decimales (valor verdadero $2.7182818\dots$).

\textbf{10.} $f(1 - x) = \frac{(1-x)^n x^n}{n!} = f(x)$. En
$\intoo{0}{1}$: $0 < x(1-x) \leq \frac14$, luego
$0 < f \leq \frac{1}{4^n\,n!}$.

\textbf{11.} $x^n(1-x)^n = \sum_{j=0}^{n} (-1)^j\binom
nj\,x^{n+j}$, luego $f = \frac{1}{n!}\sum_j c_j\,x^{n+j}$ con
$c_j \in \Z$. Por tanto, $f^{(k)}(0) = 0$ para $k < n$ o $k > 2n$, y
para $n \leq k \leq 2n$: $f^{(k)}(0) = \frac{k!}{n!}\, c_{k-n}$, un
entero porque $n! \mid k!$. La simetría da
$f^{(k)}(1) = (-1)^k f^{(k)}(0) \in \Z$.

\textbf{12.} $b^n \pi^{2n-2k} = b^n\bigl(\frac
ab\bigr)^{\,n-k} = a^{\,n-k}\,b^{\,k} \in \Z$, luego
$G(0) = \sum_k (-1)^k a^{n-k} b^k f^{(2k)}(0)$ y, análogamente,
$G(1)$ son enteros por la pregunta 11.

\textbf{13.} En $\pi^2 G + G''$, el término $k$ de $\pi^2 G$ lleva
$\pi^{2n-2k+2}f^{(2k)}$ y el término $j = k - 1$ de $G''$ lleva
$(-1)^{k-1}\pi^{2n-2k+2}f^{(2k)}$: todo se cancela salvo $k = 0$ en
la primera suma y $j = n$ en la segunda, es decir,
\[
G'' + \pi^2 G = b^n\bigl(\pi^{2n+2} f + (-1)^n f^{(2n+2)}\bigr)
= b^n \pi^{2n+2} f = \pi^2 a^n f
\]
($f$ tiene grado $2n$, luego $f^{(2n+2)} = 0$; y
$b^n\pi^{2n} = a^n$). Entonces
\[
\bigl(G'\sin\pi x - \pi G\cos\pi x\bigr)'
= (G'' + \pi^2 G)\sin \pi x = \pi^2 a^n f\sin\pi x .
\]

\textbf{14.} Integrando en $\intcc{0}{1}$:
\[
\pi^2 a^n \int_0^1 f\sin(\pi x)\,\dd x
= \bigl[G'\sin\pi x - \pi G\cos\pi x\bigr]_0^1
= \pi\bigl(G(1) + G(0)\bigr) ,
\]
luego $\pi a^n \int_0^1 f\sin\pi x = G(0) + G(1) \in \Z$. En
$\intoo{0}{1}$, $f > 0$ y $\sin\pi x > 0$: el miembro izquierdo es
positivo, luego $G(0) + G(1) \geq 1$; y $\sin \leq 1$ junto con la
pregunta 10 lo acota por $\frac{\pi a^n}{4^n\,n!}$.

\textbf{15.} Por la pregunta 1 (con $c = \frac a4$),
$\frac{\pi a^n}{4^n n!} \to 0$: para $n$ grande es $< 1$, en
contradicción con $G(0) + G(1) \geq 1$. Luego ninguna fracción
$\frac ab$ es igual a $\pi^2$: el teorema de Legendre. Y si $\pi$
fuera racional, $\pi^2$ también lo sería: así que $\pi \notin \Q$
--- y la implicación solo va en este sentido ($\sqrt2$ es irracional
con cuadrado racional), y por eso $\pi^2 \notin \Q$ es
estrictamente más fuerte que el~\cref{exo:b1:integration:9}.

\textbf{16.} Integralidad: preguntas 11--12 (\omterm{def:b1:derivative:def}{derivadas} en los
extremos); pequeñez: preguntas 10 y 1; puente: preguntas 13--14 (la
doble \omterm{thm:b1:integration:parts}{integración por partes} telescopada). Sin $\frac{1}{n!}$, los
datos en los extremos siguen siendo enteros (incluso más
fácilmente), pero la cota pasa a ser $\frac{\pi a^n}{4^n}$, que
tiende a $0$ solo cuando $a < 4$ --- y todo candidato tiene
$a = b\,\pi^2 > 9$. El factorial es exactamente lo que corre más que
el crecimiento geométrico $a^n$: sin factorial, no hay teorema.

\textbf{17.} Para $a \leq 10$, el entero $G(0) + G(1)$ es positivo y
a lo sumo $\pi\,(10/4)^n/n!$. En $n = 7$: $2.5^7 = 610.35\dots$,
luego la cota es $\frac{\pi \times 610.35}{5040} \approx 0.38 < 1$
(en $n = 6$ todavía vale $1.07$): la contradicción llega en la
séptima etapa, explícitamente.

\textbf{18.} Dos \omterm{thm:b1:integration:parts}{integraciones por partes}:
\[
\int_0^1 x(1-x)\sin\pi x\,\dd x
= \frac1\pi\int_0^1 (1 - 2x)\cos\pi x\,\dd x
= \frac{2}{\pi^2}\int_0^1 \sin\pi x\,\dd x
= \frac{2}{\pi^2}\cdot\frac{2}{\pi} = \frac{4}{\pi^3}
\]
(los términos de los corchetes se anulan: $x(1-x)$ en $0, 1$, y
$\sin\pi x$ en $0, 1$). Simbólicamente, el telescopaje de $n = 1$
(sin hipótesis sobre $\pi$) dice
$\pi^3\int_0^1 f_1\sin\pi x = -(f_1''(0) + f_1''(1))$ con
$f_1 = x(1 - x)$, $f_1'' = -2$: el miembro derecho vale $4$ --- los
dos cálculos coinciden.

\textbf{19.} $\eu^x$ resuelve $y' = y$ y $\sin\pi x$ resuelve
$y'' = -\pi^2 y$: ecuaciones lineales con coeficientes constantes,
de modo que la \omterm{thm:b1:integration:parts}{integración por partes} devuelve el núcleo sobre sí
mismo y mantiene todos los datos de \omterm{def:b1:topology:closure}{frontera} dentro de
$\Z + \Z\eu$ (resp.\ \omterm{def:b1:poly:def}{polinomios} enteros en $\pi^2$). Esa propiedad de
\omterm{def:b1:topology:closure}{clausura} es lo que la máquina necesita. Afinado con núcleos
adaptados a varios puntos a la vez, el mismo mecanismo produce el
teorema de Hermite ($\eu$ trascendente, 1873) y el de Lindemann
($\pi$ trascendente, 1882) --- fuera de este volumen.

\textbf{20.} División de \omterm{def:b1:poly:def}{polinomios} (o multiplíquese de vuelta y compruébese):
\[
x^4(1-x)^4 = (x^6 - 4x^5 + 5x^4 - 4x^2 + 4)(1 + x^2) - 4 .
\]
Integrando la identidad anterior dividida por $1 + x^2$:
\[
\int_0^1 \frac{x^4(1-x)^4}{1+x^2}\dd x
= \Bigl(\frac17 - \frac46 + 1 - \frac43 + 4\Bigr)
- 4\arctan 1
= \frac{22}{7} - \pi ,
\]
usando $\frac17 - \frac23 + 1 - \frac43 + 4 = 3 + \frac17$ y
$\arctan 1 = \frac\pi4$.

\textbf{21.} El integrando es \omterm{def:b1:continuity:continuous}{continuo} y positivo en
$\intoo{0}{1}$: la \omterm{thm:b1:integration:def}{integral} es $> 0$, luego $\pi < \frac{22}{7}$.
Además, $\frac12 \leq \frac{1}{1+x^2} \leq 1$ en $\intcc{0}{1}$ y
$\int_0^1 (x(1-x))^4 = \frac{(4!)^2}{9!} = \frac{1}{630}$
(pregunta 3):
\[
\frac{1}{1260} \leq \frac{22}{7} - \pi \leq \frac{1}{630}
\quad\Longrightarrow\quad
3.14126 < \frac{22}{7} - \frac{1}{630} \leq \pi \leq
\frac{22}{7} - \frac{1}{1260} < 3.14207 .
\]

\textbf{22.} \omterm{def:b1:complex:field}{Módulo} $x^2 + 1$: $x^2 \equiv -1$, luego
$x^{4m} = (x^2)^{2m} \equiv 1$ y
$(1 - x)^2 = 1 - 2x + x^2 \equiv -2x$, de donde
$(1-x)^{4m} \equiv (-2x)^{2m} = 4^m (x^2)^m \equiv (-4)^m$: el resto
es la constante $(-4)^m$, y el cociente $Q_m$ tiene coeficientes
enteros (división por un \omterm{def:b1:poly:def}{polinomio} entero \omterm{def:b1:poly:def}{mónico}). Dividiendo la
identidad por $1 + x^2$ e integrando:
\[
J_m := \int_0^1 \frac{(x(1-x))^{4m}}{1+x^2}\dd x
= s_m + (-4)^m\,\frac{\pi}{4},
\qquad s_m = \int_0^1 Q_m \in \Q .
\]
Despejando $\pi$: con $r_m = (-1)^{m+1}\,4^{\,1-m} s_m \in \Q$,
$\;\abs{\pi - r_m} = 4^{\,1-m} J_m \leq 4^{\,1-m}\cdot4^{-4m} =
4^{\,1-5m}$. Para $m = 1$: $s_1 = \frac{22}{7}$,
$r_1 = \frac{22}{7}$, y la cota $4^{-4} = \frac{1}{256}$ --- otra
vez las preguntas 20--21.

\textbf{23.} $\bigl|\pi - \frac{22}{7}\bigr| = \frac{22}{7} -
\pi \approx 1.26\cdot10^{-3}$, dieciséis veces mejor que la
referencia de orden $2$, $\frac{1}{7^2} \approx 2.0\cdot10^{-2}$,
que garantiza Dirichlet (\cref{pb:b1:derivative:1}, pregunta 4); y
$\bigl|\pi - \frac{355}{113}\bigr| \approx 2.7\cdot10^{-7}$ le gana
a $\frac{1}{113^2} \approx 7.8\cdot10^{-5}$ por un factor
$\approx 300$. Nada de esto contradice lo demostrado: las
desigualdades de Liouville \emph{acotan por debajo} los errores de
aproximación solo para los \omterm{pb:b1:logic:1}{números algebraicos}, y a este nivel no se
dispone de ninguna cota así para $\pi$ --- $\pi$ es libre de dejarse
aproximar espectacularmente bien.

\textbf{24.} Lema: supóngase $x = \frac pq$ y
$0 < \abs{a_n + b_n x} \to 0$. Entonces
$\abs{a_n + b_n x} = \frac{\abs{q a_n + p b_n}}{q}$, con
$q a_n + p b_n$ entero no nulo (no nulo porque el valor absoluto es
$> 0$): luego $\abs{a_n + b_n x} \geq \frac1q$ para todo $n$, en
contradicción con la convergencia a $0$. Instancias: la pregunta 7
($x = \eu$, $a_n = \alpha_n$, $b_n = \beta_n$);
el~\cref{exo:b1:seq:9} ($x = \eu$ otra vez, con
$a_q = -q!\sum_{k \leq q}\frac{1}{k!}$, $b_q = q!$); y
el~\cref{pb:b1:derivative:1}, pregunta 1, es su forma geométrica. En
la parte III y en el~\cref{exo:b1:integration:9} la trampa corre
\emph{dentro} de la contradicción: suponer la racionalidad convierte
una expresión en un entero, que el análisis aprieta después dentro de
$\intoo{0}{1}$ --- el mismo principio, traspuesto.

\textbf{25.} (i) La integralidad vive en el cálculo con
\omterm{def:b1:poly:def}{polinomios} en los extremos (preguntas 6, 11--12), la pequeñez en
las cotas que da la monotonía del $\sup$ (preguntas 4, 10) y el
puente en la \omterm{thm:b1:integration:parts}{integración por partes} (preguntas 8, 13--14) --- las
tres son teoremas de este capítulo. (ii) La \omterm{thm:b1:integration:def}{integral} aporta lo que el
teorema del valor medio no podía: una identidad \emph{exacta} entre
el objeto analítico y los datos aritméticos (una igualdad, no solo
una desigualdad con un $c$ desconocido), y por eso la máquina alcanza
$\pi^2$ mientras que el~\cref{pb:b1:derivative:1} solo alcanzaba
exponentes de aproximación. (iii) Extraído: $\eu \notin \Q$,
$\pi^2 \notin \Q$ (y por tanto $\pi \notin \Q$), $\eu = 2.718282$
certificado,
$\frac{22}{7} - \frac{1}{630} \leq \pi \leq \frac{22}{7} -
\frac{1}{1260}$ y $\binom{2n}{n} \geq \frac{4^n}{2n+1}$.
(iv) \omterm{def:b1:topology:closure}{Frontera}: Hermite y Lindemann llevan la misma máquina hasta la
trascendencia; y Apéry (1979) hizo correr la trampa del entero sobre
$\zeta(3)$ --- la máquina sigue produciendo matemáticas del siglo XX.
\end{solution}
